\chapter{FX Flow Strategies}\label{ch:s2:fx-flow-strategies}

At four in the afternoon London time the benchmark rates for currencies are fixed on the trades of a few minutes. Funds that must trade at the fix
create flows that others can anticipate, from public information about what those funds hold. Melvin and Prins found that a month in which equity
markets rise predicts currency depreciation before the month-end fix, as international equity investors rebalance their currency hedges. In public
data from 1999 to 2026, when the S\&P 500 had risen during a month, the dollar tended to fall over the month's last two days. A trade that sold it then
earned 11.4 basis points a month ($t$ = 3.4). On this chapter's synthetic market, trading ahead of an estimated hedge flow earns 12.2 basis points a
month, a Sharpe ratio of 0.99, and the return grows with the flow's estimated size. The build is \texttt{firm.fxflows}.

\section{The fix and its flows}

A benchmark fix (Book~2, chapter~17) sets a reference exchange rate from trades in a short fixing window. Asset managers value portfolios and settle
hedges at it, so a large volume of orders is executed at the fix by design. Evans studied 21 currencies over a decade. Around the 4:00 p.m.\ WMR fix,
price changes showed extraordinary volatility and negative serial correlation: prices pushed in the window tended to come back after it.

\begin{definition}[Fix flow]\label{def:s2:fx-flow-strategies:flow}
A \emph{fix flow}\index{fix flow} is the net order to buy or sell a currency at a benchmark fix that arises from mandates rather than views: hedge
rebalancing, index and portfolio rebalancing, and client orders executed at the fix; its size and direction can often be estimated in advance from
public information.
\end{definition}

The fix was also the scene of misconduct. On 12 November 2014 the Financial Conduct Authority fined five banks \pounds1\,114\,918\,000 in total for
failings in G10 spot FX between 2008 and 2013. Their traders had shared confidential information about clients' orders and attempted to manipulate
fix rates, including the 4 p.m.\ WM/Reuters and ECB fixes, and to trigger clients' stop-loss orders. Estimating a flow from public data and trading
ahead of it is legal. Trading on knowledge of a client's order, or colluding to move the fix, is not (chapter~29).

\section{Month-end rebalancing}

\begin{definition}[Rebalancing-flow estimate]\label{def:s2:fx-flow-strategies:estimate}
A \emph{rebalancing-flow estimate}\index{rebalancing-flow estimate} is a forecast of the currency that hedged international investors must buy or
sell at a rebalancing date, from the month's asset returns, estimates of their holdings, and their hedge ratios.
\end{definition}

The mechanism is arithmetic. A European investor who holds US equities and hedges half their dollar value is short dollars forward. If US equities
rise 5\% in the month, the holding is worth more and the hedge covers less than half. At month-end the investor sells more dollars forward to restore
the ratio. Summed over investors, a rising US market means dollar sales at the month-end fix. A US investor hedging foreign equities does the reverse.

The public data show the pattern. Build a dollar index from six currencies' daily noon New York rates (euro, sterling, yen, Swiss franc, Australian and
Canadian dollars), and take the S\&P 500's return from the start of each month to three days before its end. Over 332 months from 1999 to 2026, the
dollar's change over the last two days correlates with that return at $-0.18$ ($t$ = $-3.3$). Selling the dollar over those two days after a rising
month, and buying it after a falling one, earned 11.4 basis points a month ($t$ = 3.4)
(\cref{fig:s2:fx-flow-strategies:real}). Over the last three days the trade earned 7.4 ($t$ = 1.9), and over the last five 13.7 ($t$ = 2.8). The rates
are taken at noon New York rather than at the London fix, so the test sees the pattern at one remove.

\begin{omfigure}
\begin{tikzpicture}
\begin{axis}[width=8cm,height=5.2cm,ybar,bar width=22pt,xlabel={\small last $k$ trading days of the month},ylabel={\small trade (bp a month)},
  tick label style={font=\footnotesize},xmin=1,xmax=6,ymin=0,ymax=16,xtick={2,3,5},grid=major,grid style={gray!25},table/col sep=comma,
  nodes near coords,every node near coord/.append style={font=\scriptsize}]
\addplot[fill=omThm!70,draw=omThm] table[x=k,y=bp]{figdata/strategies-2/15-fx-flow-strategies/real.csv};
\end{axis}
\end{tikzpicture}

\omcaption{A month-end hedge-flow trade in public data: sell a six-currency dollar index over the month's last $k$ days when the S\&P 500 has risen
since the month began, buy it when the S\&P has fallen; mean return a month, January 1999 to September 2026 (332 months). From FRED's H.10 rates and
Cboe's S\&P 500 history; derived statistics only. Data: \texttt{s2\_fetch\_monthend}.}\label{fig:s2:fx-flow-strategies:real}
\end{omfigure}

\texttt{firm.fxflows} builds the mechanism explicitly (\cref{lst:s2:fx-flow-strategies:model}). Over 360 synthetic months, foreign investors hold home
equities hedged at a ratio of 0.5, and home investors hold foreign equities hedged at 0.3. The month-end flow moves the exchange rate in the last days
before the fix by 12.2 basis points in standard deviation, and half of the move reverses after. The trader estimates the flow with 30\% errors in the
holdings and hedge ratios; the estimate still correlates with the truth at 0.94. Trading the estimated flow's direction earned 12.2 basis points a
month after costs ($t$ = 5.4), a Sharpe ratio of 0.99. Fading the move in the days after month-end earned 10.2 ($t$ = 4.9). By fifth of estimated flow
size, the trade earned 0.2, 12.0, 17.7, 9.5 and 21.8 basis points against expected pushes of 1.3, 4.5, 7.7, 12.8 and 23.4
(\cref{fig:s2:fx-flow-strategies:quintiles}).

\begin{omfigure}
\begin{tikzpicture}
\begin{axis}[width=10cm,height=5.4cm,ybar,bar width=10pt,xlabel={\small fifth of months by estimated flow size},ylabel={\small bp a month},
  tick label style={font=\footnotesize},xmin=0.4,xmax=5.6,ymin=0,ymax=26,xtick={1,2,3,4,5},grid=major,grid style={gray!25},table/col sep=comma,
  legend style={font=\scriptsize,at={(0.5,-0.3)},anchor=north},legend columns=2]
\addplot[fill=gray!40,draw=gray] table[x=q,y=est]{figdata/strategies-2/15-fx-flow-strategies/quintiles.csv};
\addplot[fill=omThm!70,draw=omThm] table[x=q,y=realised]{figdata/strategies-2/15-fx-flow-strategies/quintiles.csv};
\legend{expected push from the estimate,trade's mean return}
\end{axis}
\end{tikzpicture}

\omcaption{The synthetic month-end flow trade by fifth of the estimated flow's size: the push the estimate implies and the trade's mean return after
costs, 360 months. Data: \texttt{s2\_fxflows.results}.}\label{fig:s2:fx-flow-strategies:quintiles}
\end{omfigure}

\section{Central-bank behaviour}

\begin{definition}[Central-bank reaction trade]\label{def:s2:fx-flow-strategies:cb}
A \emph{central-bank reaction trade}\index{central-bank reaction trade} positions for a central bank's predictable response to its currency, such as
intervention near levels it has defended before or reserve rebalancing after large moves, and fades or follows the move depending on the bank's
record and its capacity to act.
\end{definition}

Central banks are the largest flow traders there are, and some are predictable. A bank that has intervened to stop a depreciation beyond some level
may do so again. A reserve manager that rebalances its currency shares after a large move creates flows like a hedger's. The trades are rare and
the counterparty can print money. A trader who fades a central bank's stated line is betting on its will to act and its resources to do so. The chapter
gives no figures for this trade: intervention records are published late and in varied forms, and none is used here.

\section{Trading predictable flows}

In the fix window itself the synthetic flow pushes the rate by 30\% of its month-end move, and 40\% of that push reverses within the next half hour.
A liquidity provider who takes the other side of the push earned 0.75 basis points a month after a 0.5 basis-point cost ($t$ = 2.7), a Sharpe ratio of
0.50. That is the negative serial correlation Evans found, seen from the side that supplies liquidity rather than the side that demands it.

Three conditions make a flow trade legitimate and durable. The flow must be estimated from public information. It must be large relative to the
market's depth over the window. And the trader must be willing to be wrong on the months when the estimate misses, because the estimate is noisy and
other traders may be on the same side.

\section{Strategy files}

\begin{strategyfile}{Month-end hedge-rebalancing flow}\label{strat:s2:fx-flow-strategies:monthend}
\sfield{Who pays you, and why} International equity and bond investors who rebalance their currency hedges on a fixed schedule, whatever the
price.
\sfield{Instruments and venues} Spot and forwards in the major pairs, into the month-end fix.
\sfield{Signal} The month's equity (and bond) returns by country, times estimated holdings and hedge ratios.
\sfield{Sizing and execution} In proportion to the estimated flow; out at or after the fix.
\sfield{Costs} Spreads near the fix, when they widen.
\sfield{How it dies} Hedgers move their rebalancing off the fix; too many traders anticipate the same flow.
\sfield{Horizon, capacity, infrastructure} Days; flow estimates by country and asset.
\sfield{Backtest honestly} Holdings and hedge ratios as they could be known then; fix-time prices.
\sfield{Sources} Melvin and Prins (2015); public-data test: 11.4 basis points a month over the last two days, 1999--2026; this chapter: 12.2 in the
synthetic market.
\end{strategyfile}

\begin{strategyfile}{Fix-window liquidity provision}\label{strat:s2:fx-flow-strategies:fix}
\sfield{Who pays you, and why} Orders that must execute at the fix and push the rate in the window.
\sfield{Instruments and venues} Spot in the fixing window.
\sfield{Signal} The window's order imbalance as it forms.
\sfield{Sizing and execution} Take the other side in the window; close within the hour.
\sfield{Costs} Window spreads.
\sfield{How it dies} Fix methodology changes that spread orders over a longer window.
\sfield{Horizon, capacity, infrastructure} Minutes; low-latency spot access.
\sfield{Backtest honestly} Tick data in the window; no knowledge of clients' orders.
\sfield{Sources} Evans (2018) on reversal around the fix; this chapter: 0.75 basis points a month.
\end{strategyfile}

\begin{strategyfile}{Central-bank intervention fade}\label{strat:s2:fx-flow-strategies:cb}
\sfield{Who pays you, and why} Speculators who push a currency against a central bank's defended line and must retreat when it intervenes.
\sfield{Instruments and venues} Spot and options in the pair.
\sfield{Signal} Distance to past intervention levels; official statements.
\sfield{Sizing and execution} Small, with options to bound the loss if the line breaks.
\sfield{Costs} Option premiums.
\sfield{How it dies} The bank gives up the line.
\sfield{Horizon, capacity, infrastructure} Days to weeks.
\sfield{Backtest honestly} Interventions as published, often weeks later.
\sfield{Sources} No performance figure verified.
\end{strategyfile}

\begin{strategyfile}{Equity-performance FX hedge flow}\label{strat:s2:fx-flow-strategies:equityhedge}
\sfield{Who pays you, and why} As for the month-end flow, measured from the relative performance of two equity markets.
\sfield{Instruments and venues} The pair between the two markets' currencies.
\sfield{Signal} Home against foreign equity performance over the month.
\sfield{Sizing and execution} Last days before month-end.
\sfield{Costs} Spreads.
\sfield{How it dies} Hedge ratios change; other flows dominate.
\sfield{Horizon, capacity, infrastructure} Days.
\sfield{Backtest honestly} Both markets' closes before the trade.
\sfield{Sources} Melvin and Prins (2015); no performance figure verified beyond this chapter's tests.
\end{strategyfile}

\section{Tutorial: four o'clock}\label{tut:s2:fx-flow-strategies:tut}

\begin{tutorial}
\textbf{Goal.} Simulate month-end hedge flows, estimate them, trade ahead of them and provide liquidity at the fix; test the pattern on public data.
\textbf{End state:} the numbers in the text and the two figures.
\begin{tutsteps}
\item \textbf{Flows and trades}.
\omcode{code/firm/fxflows/firm_fxflows.py}{46}{76}{Month-end flows from equity returns and hedge ratios, and the two trades.}\label{lst:s2:fx-flow-strategies:model}
\item \textbf{Results}.
\omcode{code/strategies-2/15-fx-flow-strategies/python/s2_fxflows.py}{39}{51}{The flow trade by estimated size, the fade and the fix trade.}\label{lst:s2:fx-flow-strategies:results}
\item \textbf{Run} \texttt{s2\_fetch\_monthend.py} once, then \texttt{results()} and \texttt{fig\_fxflows.py}.
\end{tutsteps}
\textbf{What to change next.} Let hedge ratios drift over the years; add a second, uncorrelated flow (index rebalancing) that sometimes offsets the
first; trade only the two fifths with the largest estimated flow.
\end{tutorial}

\section{Build: FX flows}\label{bld:s2:fx-flow-strategies:fxflows}

\begin{build}
\bfield{Purpose} Month-end hedge flows from asset performance and hedge ratios, a fix-window push and reversal, and the trades around them.
\bfield{Interface} \texttt{FlowFXConfig(\dots)}, \texttt{simulate\_months(cfg)}, \texttt{flow\_trade(sim, cfg, scale)},
\texttt{fix\_liquidity(sim, cfg)}.
\bfield{Rules} The flow's sign follows the hedges; the estimate uses noisy holdings and hedge ratios; costs per trade.
\bfield{Acceptance tests} \texttt{code/firm/fxflows/tests/}: the flow's sign and size by hand; without noise the trades earn the push, or its
reversal, less costs.
\bfield{Stretch} Several currencies; bond hedges; index rebalancing.
\end{build}

\begin{omsources}
\item M.~Melvin and J.~Prins, ``Equity hedging and exchange rates at the London 4 p.m.\ fix'', \emph{Journal of Financial Markets} 22, 2015.
\item M.~D.~D.~Evans, ``Forex trading and the WMR Fix'', \emph{Journal of Banking and Finance} 87, 2018.
\item Financial Conduct Authority, press release of 12 November 2014 on FX failings at five banks.
\item Board of Governors H.10 exchange rates, via FRED; Cboe S\&P 500 index history.
\end{omsources}

\section{Exercises}

\begin{exercise}[$\star$]\label{exo:s2:fx-flow-strategies:1}
A European investor holds \$100 million of US equities hedged at 50\%. US equities rise 5\% in the month. How many dollars must it sell forward at
month-end?
\end{exercise}

\begin{exercise}[$\star$]\label{exo:s2:fx-flow-strategies:2}
A trade earns 11.4 basis points a month with $t$ = 3.4 over 332 months. What is its annual Sharpe ratio?
\end{exercise}

\begin{exercise}[$\star$]\label{exo:s2:fx-flow-strategies:3}
What was the average FCA fine per bank in November 2014?
\end{exercise}

\begin{exercise}[$\star\star$]\label{exo:s2:fx-flow-strategies:4}
Why does a rising US equity market lead to dollar sales at month-end?
\end{exercise}

\begin{exercise}[$\star\star$]\label{exo:s2:fx-flow-strategies:5}
Why is the public-data test only an approximation of the fix trade?
\end{exercise}

\begin{exercise}[$\star\star$]\label{exo:s2:fx-flow-strategies:6}
What separates a legitimate flow trade from the conduct the FCA fined?
\end{exercise}

\begin{exercise}[$\star\star\star$]\label{exo:s2:fx-flow-strategies:7}
\emph{Coding.} Rerun with \texttt{FlowFXConfig(est\_error=1.0)}. How do the estimate's correlation with the true flow and the trade's return change?
\end{exercise}

\begin{exercise}[$\star\star\star$]\label{exo:s2:fx-flow-strategies:8}
\emph{Find the flaw.} ``The flow trade works in every fifth of estimated size, so we should size it equally every month.''
\end{exercise}

\section{Problem: Four O'Clock}

\begin{problem}[{Weekend problem --- predictable FX flows}]\label{pb:s2:fx-flow-strategies:1}
The chapter's synthetic flows, public data and the public record.

\medskip\noindent\textbf{Part I --- The fix.}
\begin{enumerate}
\item What is a benchmark fix, and why is so much traded at it?
\item What did Evans find around the fix?
\item What did the FCA find in 2014?
\item Define a \omterm{def:s2:fx-flow-strategies:flow}{fix flow}.
\end{enumerate}

\medskip\noindent\textbf{Part II --- Month-end.}
\begin{enumerate}[resume]
\item Explain the hedge-rebalancing flow.
\item What did Melvin and Prins find?
\item Give the public-data test's results.
\item Define a \omterm{def:s2:fx-flow-strategies:estimate}{rebalancing-flow estimate}.
\end{enumerate}

\medskip\noindent\textbf{Part III --- The synthetic market.}
\begin{enumerate}[resume]
\item Describe the synthetic flows and the trader's estimate.
\item Give the flow trade's return and Sharpe ratio.
\item How does the return vary with the estimated size?
\item What does the liquidity provider earn in the fix window?
\end{enumerate}

\medskip\noindent\textbf{Part IV --- The verdict.}
\begin{enumerate}[resume]
\item State the \emph{named result}: the month-end flow trade's return against the estimated flow size.
\item What does fading after month-end earn?
\item Define a \omterm{def:s2:fx-flow-strategies:cb}{central-bank reaction trade} and its risk.
\item What makes a flow trade legitimate?
\item How would you backtest the month-end trade honestly?
\item Which strategy file needs the fastest execution?
\item How does this chapter relate to chapter~5's rebalancing flows?
\item In one sentence: who pays the flow trader?
\end{enumerate}
\end{problem}

\section{Interview questions}

\begin{interviewq}[$\star$ \iqroles{trader}]\label{iq:s2:fx-flow-strategies:1}
What is the WM/Reuters 4 p.m.\ fix, and why does it matter?
\end{interviewq}

\begin{interviewq}[$\star\star$ \iqroles{researcher}]\label{iq:s2:fx-flow-strategies:2}
How would you estimate month-end FX hedge flows?
\end{interviewq}

\begin{interviewq}[$\star\star$ \iqroles{trader}]\label{iq:s2:fx-flow-strategies:3}
A client asks you to buy \$2 billion at the fix. What may you do before the fix, and what may you not?
\end{interviewq}

\begin{interviewq}[$\star\star$ \iqroles{risk}]\label{iq:s2:fx-flow-strategies:4}
What are the risks of a month-end flow book?
\end{interviewq}

\begin{interviewq}[$\star\star$ \iqroles{developer}]\label{iq:s2:fx-flow-strategies:5}
Design a monitor for trading around benchmark fixes that a compliance team would accept.
\end{interviewq}

\begin{interviewq}[$\star\star\star$ \iqroles{researcher}]\label{iq:s2:fx-flow-strategies:6}
With holdings $H$, hedge ratio $h$ and an asset return $R$, show that restoring the hedge requires a currency trade of $hHR$, and find the estimate's
error when $H$ and $h$ are each off by a factor $(1 + \epsilon)$.
\end{interviewq}
